\section{Related Work}
\label{sec:related}

\subsection{Multi-armed and contextual bandits}

The multi-armed bandit problem dates to Thompson \cite{thompson1933likelihood}.
The modern era is anchored by Auer et al.'s UCB1 and its variants
\cite{auer2002finite,audibert2009exploration}, by Russo and Van Roy's
analysis of Thompson Sampling
\cite{russo2014learning,russo2018tutorial,agrawal2012analysis}, and by Bubeck
and Cesa-Bianchi's regret theory \cite{bubeck2012regret}.

Contextual extensions begin with Auer's LinRel and progress through
Li et al.'s LinUCB \cite{li2010contextual}, Agrawal and Goyal's contextual
Thompson sampling \cite{agrawal2013thompson}, and the neural-bandit family
\cite{zhou2020neural,zhang2020neural}. Microsoft Research's open-source
Decision Service \cite{agarwal2016making} exemplifies a production-grade
contextual-bandit deployment with off-policy evaluation built in.

\subsection{Constrained and conservative bandits}

The conservative bandit setting --- where the learner must perform no
worse than a known baseline with high probability --- was introduced by Wu
et al.\ \cite{wu2016conservative} and extended to linear payoffs by Amani
et al.\ \cite{amani2019linear}. Kazerouni et al.\ \cite{kazerouni2017conservative}
analyse the related setting under linear contextual rewards. These works
provide the formal scaffolding for our conservative-property guarantee:
the quality-preservation constraint is structurally identical to a
no-worse-than-baseline constraint with the opus tier playing the baseline
role.

Safe exploration in reinforcement learning
\cite{garcia2015comprehensive,achiam2017constrained} addresses adjacent
concerns; we do not adopt the full RL apparatus because our decisions are
single-step.

\subsection{Calibration and uncertainty quantification}

Platt scaling \cite{platt1999probabilistic} and isotonic regression
\cite{zadrozny2002transforming} are the classical methods for converting
classifier scores into calibrated probabilities. We use them only as
diagnostic checks; the operational calibration in CC-TS is provided by
conformal prediction \cite{vovk2005algorithmic,shafer2008tutorial,angelopoulos2021gentle},
which yields distribution-free coverage guarantees on the
quality-drop intervals.

Hierarchical Bayesian modelling with partial pooling
\cite{gelman2013bayesian} underpins our prior structure. The prior at the
agent--shape level is partially pooled with both the agent-level marginal
and the shape-level marginal, which dramatically improves cold-start
behaviour for low-traffic cells.

\subsection{LLM model routing and cost optimization}

A recent line of work attacks the model-routing problem directly.
RouteLLM \cite{ong2024routellm} learns a binary router between strong and
weak models; FrugalGPT \cite{chen2023frugalgpt} composes models in
escalation cascades. Hybrid-LLM \cite{ding2024hybrid} treats routing as a
contextual bandit. Lu et al.\ \cite{lu2024routing} compare against an
oracle. None of this work, to our knowledge, embeds an asymmetric loss
matrix or a hard quality-preservation constraint --- their decision rules
are optimised for accuracy on a benchmark, not for not-regressing on a
production workload.

\subsection{Agent frameworks}

\blackrim sits in a small but growing class of multi-agent frameworks
oriented around code: AutoGen \cite{wu2023autogen}, MetaGPT
\cite{hong2023metagpt}, ChatDev \cite{qian2023communicative}, and
Anthropic's Claude Code itself
\cite{anthropic2024claudecode}. Routing decisions in these frameworks are
either agent-static (\blackrim's pre-advisor state) or based on simple
keyword heuristics. The advisor we describe is, to our knowledge, the
first to bring constrained Bayesian decision theory to multi-agent LLM
routing.

\subsection{Off-policy evaluation}

Inverse propensity scoring \cite{horvitz1952generalization,dudik2011doubly}
and doubly-robust estimators
\cite{dudik2014doubly,jiang2016doubly,thomas2016data} provide the
theoretical machinery for evaluating a counterfactual policy from logged
data. We use these estimators in \cref{sec:eval} to estimate what each
baseline would have spent on the same dispatch trace.
